\section{Introduction}
\label{sec:intro}

% Opening paragraph: the three levels of the RSI loop, and why the algorithmic one
% is the one that compounds. The claim is deliberately *not* that compute and data
% saturate -- a power law has no ceiling, it only gets more expensive. The claim is
% that the three levels differ in whether their result is inherited by the next run.
Recursive self-improvement (RSI) rests on a loop: a system improves the process that produces its successor, and the successor inherits the improvement. 
Three levels in that loop can be automated by a coding agent --- the \emph{Systems engineering} level (kernels, parallelism, communication), the \emph{Data} level (mixtures, synthesis, filtering), and the \emph{Algorithmic design} level (objectives, update rules, regularization,
schedules). 
Work on the \emph{Systems engineering} level is bounded by the hardware it runs on, since once a kernel reaches the hardware roofline, further gains are no longer possible.
Work on the \emph{Data} level is bounded by a finite stock of human text \citep{villalobos2022rundata}, by synthetic corpora that
largely re-express what the model already carries \citep{shumailov2023curse}, and by a power law in which every further halving of the loss gap costs several times more tokens than the halving before it \citep{kaplan2020scaling,hoffmann2022chinchilla}. 
Work on the \emph{Algorithmic design} level is different in kind: a better objective or update rule changes the exchange rate between compute and capability \citep{ho2024algprogress}, so every training run that follows --- including the run that produces the next agent --- inherits the gain. Adam, layer normalization, DPO and GRPO were each paid for once and have been earning since; if RSI is going to compound, most of the compounding has to come from this level.

% Paragraph 2: no existing benchmark isolates the algorithmic level. Each sentence
% names what a benchmark's *contract* rewards, rather than grading its quality.
No benchmark isolates that level. The Kaggle-derived suites ask for a competition submission ---predictions over one fixed dataset --- so what wins is feature engineering and ensembling while the
learning algorithm stays a library call the agent never edits, and nothing in the submission is
inherited by the next training run \citep{chan2024mlebench,qiang2025mledojo}. PostTrainBench states
its task end to end, post-training a base model against a released instruction-tuned checkpoint, and
its largest levers are which data to assemble and what to initialize from rather than the objective
\citep{rank2026posttrainbench}; RSIBench-Data makes that emphasis deliberate, freezing the
post-training stack so that only data-centric decisions vary \citep{meng2026rsibench}. MLS-Bench
comes closest, with 140 tasks in which an agent improves one component of an ML system
\citep{lyu2026mlsbench}, but the component boundary is handed to it and the score conflates execution-level improvements with changes to the learning algorithm.
Closest of all is
\texttt{autoresearch}, which hands the agent one training file and declares everything in it fair
game, architecture and optimizer included \citep{karpathy2026autoresearch}; but a five-minute run of a single script is not a research repository, and measured against classical optimizers the agent's
edits behave like hyperparameter search and lose to CMA-ES and TPE \citep{ferreira2026autoresearch}. Editing source code is therefore not
the same as designing an algorithm, and none of these settings answers the question the loop turns
on: did the agent change \emph{how the run was executed}, or \emph{how the model learns}?

% Paragraph 3: the operational distinction. This is the definition the patch
% taxonomy in Section 5 implements, so it has to be stated mechanically here.
The line between the two is not the size of a change but what it touches: a hyperparameter is a number
the training algorithm takes as given, an algorithmic change rewrites the algorithm --- the loss it
optimizes, the update it applies. 
The second kind is what a machine learning scientist does at an industrial training system. They read the training dynamics --- the loss curve and where it spikes, gradient norms,
the entropy of the policy, the divergence from the reference model, the distribution of advantages,
the loss broken down by token --- and infer from them which part of the algorithm is misbehaving: a
policy whose entropy has collapsed, a penalty term that has come to dominate the objective, a reward
model saturating on the easy half of its pairs. The diagnosis names a mechanism and the fix changes
that mechanism.

% Paragraph 4: the apparatus. Terminology fixed here and to be kept everywhere:
%   - "training algorithm" for the object of improvement (never "recipe", which in
%     earlier drafts meant the committed configuration, that same configuration at
%     a longer budget, and the method in general, all in one word)
%   - "baseline" for whatever a score is compared against, always named
%   - "system" = (model, harness, reasoning effort); "cell" = (system, task)
%   - patches are "classified", not "audited"; the corpus is a grid, not a census
In this paper, we propose \textbf{\ai{}}, a benchmark built to isolate the algorithmic design level. It includes 10
research repositories, each representing a distinct family of training algorithms --- supervised fine-tuning,
multi-turn agentic RL, on-policy distillation, Bradley--Terry reward modeling, preference
optimization, diffusion RL, machine unlearning, discrete graph diffusion, weight averaging and
one-shot pruning --- and asks an agent to improve each repository's \emph{own} training algorithm,
rather than to reach a target somebody else set. Every task carries the same contract. The agent
has 4 hours on one B300 GPU to read the repository, change its training code, and test each idea
against a fast proxy metric. When the 4 hours are up the agent stops, and the code it leaves
behind is trained from a clean start for up to 12 hours. This is the asymmetry a machine learning
scientist works under: an idea can be triaged in minutes, but the run that settles it takes a day.
The data behind that final measurement is never available during the 4 hours: the agent may
consult its proxy as often as it likes, but the evaluation that decides its score is out of reach,
exactly as a held-out test set is out of reach of the development loop in any industrial training
system.

% Paragraph 5: what the agents submit, in three steps that build on each other.
% Numbers computed from build/family_rows.jsonl joined to the final experiment
% sheet. NPO is a single balance score with baseline 0.95201, so it is included.
\textbf{Results.} Across 29 configurations of six systems on all ten tasks the mean
score is $0.166$ and the best system reaches $0.250$, on a scale where $0.1$ is the
algorithm the repository ships and $1.0$ the task optimum. The submissions say
where the remaining distance went. Of the 263 that change anything, 141 leave the
learning procedure exactly as they found it and move budgets, checkpointing,
hyperparameters and capacity instead. The 122 that do reach it --- the objective,
the supervision signal, the learning rule, the data --- average $0.226$ against
$0.126$ for the rest: the algorithmic layer is where the distance gets closed, and
most submissions never go there. More reasoning effort moves agents towards it,
taking that share from $8\%$ to $64\%$ and the mean score from $0.094$ to $0.196$,
which is most of what effort buys. A machine learning scientist reads the training
dynamics, names the mechanism that is failing, and changes that mechanism; few of
these submissions do.

% Paragraph 6: contributions. Deliberately three, not 4: the baseline
% sensitivity result is real but stays out until the sheet's Baseline row has a
% stated definition.
In summary, we make the following contributions:
\begin{itemize}[leftmargin=1.2em,itemsep=3pt,topsep=3pt,parsep=0pt]
  \item \textbf{A benchmark that isolates the algorithmic level.} Ten research repositories spanning ten families of training algorithms, evaluated under a unified protocol that separates a four-hour agent development window from the up to twelve-hour clean-start training run used for scoring.
  \item \textbf{A measurement of what agents \emph{do}, not only whether they win.} Every submission is classified by what it changes, which is what turns
    ``the agent improved the training algorithm'' into a checkable statement
    rather than a restatement of the score.
\item \textbf{Revealing a gap between algorithmic exploration and actual improvement.}
    The algorithmic design level is where the distance to a better algorithm is closed and the one these agents reach least often, and more reasoning effort mostly buys the willingness to reach it.
\end{itemize}
% We release every artifact and record behind these numbers.

% TODO: remaining introduction paragraphs.
% (4a) parked from paragraph 4: the validity argument for splitting the budgets --
%      had the agent simply been given twelve hours it would have been steering the
%      run we report, and a good score could mean a better training algorithm or an
%      agent that stopped a run at a lucky moment. Belongs with the apparatus.
% (4b) the experiment grid -- six systems, 29 (model, effort) configurations, 290
%      configuration-task cells, 271 of them scored -- and the fact that every
%      submitted patch is classified by what it touches. Cut from paragraph 4 so
%      that it states the contract only; both belong with the results.
% (5) headline result + the run-management/hyperparameter vs learning-rule split.
%     Blocked: the final experiment sheet's Baseline row needs a definition, and
%     resolved: NPO is a balance score, higher better, baseline 0.95201.
% (6) baseline choice as the second contribution: the same patches move from 0%
%     to 82% on OpenR1 and from 100% to 44% on DiGress when the baseline changes.
% Parked from earlier drafts of this section, for use later rather than here:
%   - the classification criterion in full (which term enters the loss, how the
%     update is formed, what is regularized, which gradient is applied and when)
%     belongs in the taxonomy section, not the introduction
%   - "a one-line diff can cross it and a thousand-line refactor need not, which
%     is why it has to be read off the patch rather than off the score" belongs
%     in paragraph 4's account of why patches are extracted and classified
